\begin{frame}[plain]
    \begin{center}
        \LARGE Pré-processamento de imagens
    \end{center}
\end{frame}

\begin{frame}{Pré-processamento}
    O pré-processamento tenta simular os ruídos e alterações na imagem que pode ocorrer durante a sua captura, armazenamento ou transmissão. 
\end{frame}

% \begin{frame}{Brilho}
    
% \end{frame}

\begin{frame}[plain]
\begin{figure}
    \centering
    \caption{Imagem com $0.5$ de brilho}
    \includegraphics[width=\textwidth]{images/pre-processing/Low brightness 0.5.png}
    \begin{center}
        {\footnotesize Fonte: Chevanon Photography, 2018; modificação aplicado pelo autor.}
    \end{center}    
    \label{fig:enter-label}
\end{figure}
\end{frame}

\begin{frame}[plain]
    \begin{figure}
        \centering
        \caption{Imagem com $1.5$ de brilho}
        \includegraphics[width=\textwidth]{images/pre-processing/High brightness 1.5.png}
        \begin{center}
        {\footnotesize Fonte: Chevanon Photography, 2018; modificação aplicado pelo autor.}
    \end{center}
        \label{fig:enter-label}
    \end{figure}
\end{frame}

% \begin{frame}{Contraste}
    
% \end{frame}

\begin{frame}[plain]
    \begin{figure}
        \centering
        \caption{Imagem com $0.5$ de contraste}
        \includegraphics[width=\textwidth]{images/pre-processing/Low contrast 0.5.png}
        \begin{center}
        {\footnotesize Fonte: Chevanon Photography, 2018; modificação aplicado pelo autor.}
    \end{center}
        \label{fig:enter-label}
    \end{figure}
\end{frame}

\begin{frame}[plain]
    \begin{figure}
        \centering
        \caption{Imagem com $1.5$ de contraste}
        \includegraphics[width=\textwidth]{images/pre-processing/High contrast 1.5.png}
        \begin{center}
        {\footnotesize Fonte: Chevanon Photography, 2018; modificação aplicado pelo autor.}
    \end{center}
        \label{fig:enter-label}
    \end{figure}
\end{frame}

% \begin{frame}{\textit{Gaussian blur}}
    
% \end{frame}

\begin{frame}[plain]
    \begin{figure}
        \centering
        \caption{Imagem com \textit{Gaussian blur} com desvio padrão de $101$}
        \includegraphics[width=\textwidth]{images/pre-processing/Gaussian blur 101.png}
        \begin{center}
        {\footnotesize Fonte: Chevanon Photography, 2018; ruído aplicado pelo autor.}
    \end{center}
        \label{fig:enter-label}
    \end{figure}
\end{frame}

% \begin{frame}{\textit{Gaussian noise}}
    
% \end{frame}

\begin{frame}[plain]
    \begin{figure}
        \centering
        \caption{Imagem com \textit{Gaussian noise} com desvio padrão de 90}
        \includegraphics[width=\textwidth]{images/pre-processing/Gaussian noise 90.png}
        \begin{center}
        {\footnotesize Fonte: Chevanon Photography, 2018; ruído aplicado pelo autor.}
    \end{center}
        \label{fig:enter-label}
    \end{figure}
\end{frame}

% \begin{frame}{\textit{Box blur}}
    
% \end{frame}

\begin{frame}[plain]
    \begin{figure}
        \centering
        \caption{Imagem com \textit{box blur} com 51 \textit{pixels} para média}
        \includegraphics[width=\textwidth]{images/pre-processing/Box blur 51.png}
        \begin{center}
        {\footnotesize Fonte: Chevanon Photography, 2018; ruído aplicado pelo autor.}
    \end{center}
        \label{fig:enter-label}
    \end{figure}
\end{frame}

% \begin{frame}{Sal e pimenta (\textit{Salt and pepper})}
    
% \end{frame}

\begin{frame}[plain]
    \begin{figure}
        \centering
        \caption{Imagem com sal e pimenta com probabilidade de $0.5$ de ocorrência}
        \includegraphics[width=\textwidth]{images/pre-processing/Salt and pepper 0.5.png}
        \begin{center}
        {\footnotesize Fonte: Chevanon Photography, 2018; ruído aplicado pelo autor.}
    \end{center}
        \label{fig:enter-label}
    \end{figure}
\end{frame}